\section{Discussion}

Existing selected-data studies indicate that the declared Sample-I and Sample-II populations occupy the B stack differently in depth. The difference has now been carried through the pre-threshold 8-by-16 event reconstruction with run-block uncertainty and a baseline/threshold scan, providing the paper's most direct experimental range-stack observation: its direction is stable across the 0--1000~ADC scan while its magnitude is strongly threshold-compressed. Even so, the first interpretation remains a run-family dependence: a unique trigger-causation statement additionally requires the primary hardware trigger and run log to be source-bound.

Sparse longitudinal sampling is a central detector-design issue. An alternating four-position readout can preserve a coarse penetration pattern while losing the local sequence of energy loss near the Bragg rise. This creates a concrete aliasing mechanism: changes in particle energy, incidence angle or passive material can move the high-\(dE/dx\) region between instrumented positions. A quantitative PID or range estimator based on sparse \(\Delta E\)--\(E\) information therefore requires either sufficiently dense longitudinal instrumentation or an independently validated model of the missing-layer response, with the readout phase treated as a nuisance rather than as hidden truth.

The timing analysis gives a complementary negative result. In the component-safe rerun under the measured polarity map, the B4--B6 pair residual has a central-68\% width of 0.146~ns (bootstrap interval 0.144--0.148~ns) with a constant $-9.7$~ns channel-skew median, superseding both the earlier 38~ns direct-mode diagnostic and the 8.748~ns value produced under the falsified polarity convention. This width is a high-SNR sample-interpolation precision at fixed amplitude selection, not an intrinsic stave resolution, and it is not divided by $\sqrt{2}$. Component identity, timing reference, run dependence and channel covariance must be established before any deconvolution to a single-stave resolution is attempted. This limitation is scientifically useful because it states what a future campaign must record or calibrate: a source-bound timebase, sufficient waveform window, stable component assignment and an independent timing reference or covariance model.

The optical-model audit shows why a detector-response study must distinguish a causal response graph from a tuned total efficiency. WLS fluorescence multiplicity, trapping and attenuation, fibre-end coupling, SiPM photon-detection response, microcell occupancy/recovery, correlated noise and electronics can compensate one another in an integrated charge or amplitude. Stage-resolved counters, phase-space scans and independent bench/hardware constraints are therefore more informative than tuning a single PE/ADC conversion to one dataset.

The present limitations are coupled rather than independent. Geometry and run/trigger uncertainty affect which particles enter the selected data population; source-model and event-measure uncertainty affect the MC comparison; sparse readout affects the observability of a local stopping-power rise; and optical/electronics assumptions affect any conversion from deposited energy to charge or ADC. The final paper should therefore report the direct data observables first and introduce model-dependent detector-performance claims only after these nuisance layers have been propagated through the same final population definitions.
